% Обмеження і загрози валідності.
\section{Обмеження і загрози валідності}\label{sec:limitations}

\subsection{Корпус і запити}
\uapar{} імітує особистий архів відкритими матеріалами: статті Вікіпедії представляють нотатки, а їхні перетворені версії --- офісні документи. Тому текст однорідніший і краще відредагований за справжні нотатки, а приватних матеріалів у ньому немає. 200 запитів написала одна модель за правилами, а не люди, що шукали власні документи; усі цілі запитів із відповіддю --- нотатки. Окреме розширення з 50 запитів охоплює документи, таблиці й файли репозиторіїв (розділ~\ref{sec:results-extension}), але також написане моделлю. Воно містить 49 різних формулювань, тому 50 записів не цілком незалежні; окремо подано детермінований аналіз чутливості на унікальних формулюваннях. Запити без відповіді за побудовою не мають лексичних збігів, що наперед визначає результат лексичних систем на них. Теми відповідають інтересам автора; інший архів може сприяти іншим системам.

\subsection{Оцінки релевантності від мовних моделей}
Текстову релевантність оцінюють мовні моделі, а не незалежна група людей. Узгодженість, перевірки цільових документів і навмисно безвідповідних запитів дають діагностичні свідчення, проте не утворюють повного еталона. Gemini та Gemma можуть мати спільні помилки через спільного постачальника. Повторний перегляд Gemini не усуває цю залежність; крім того, під час нього видно початкові оцінки. Тому початкові мітки зберігаються, а метрики й порядок систем порівнюються за різними суддями та до і після перегляду. Основного суддю обрано після аналізу його діагностики --- це методичне рішення, ухвалене після отримання оцінок. Ця версія свідомо не містить людського аудиту релевантності. Тому висновки щодо текстового пошуку залежать від модельних суджень; незалежна людська перевірка є майбутньою роботою.

\subsection{Охоплення запитів і кандидатів}
Корпус неоднорідний, але всі задані цілі запитів із відповіддю в початковому наборі з 200 запитів --- нотатки. Окреме розширення додає малу вибірку документів, таблиць і файлів репозиторіїв. Жоден набір не встановлює загальної якості для довільних форматів цілей, зображень чи записів. E5 шукає в меншому текстовому наборі документів, а індекси EG2 включають медіа. Контроль на однакових кандидатах (табл.~\ref{tab:encoder-control}) змінює nDCG@10 EG2 на $-0{,}0015$ --- малу точкову величину в цьому контролі. Це не усуває всієї невизначеності щодо охоплення: контрольні прогони не формували пул, а 0,29\% документів їхньої першої десятки не оцінено.

Повністю оцінено лише об'єднання перших десяти результатів усіх систем. Поза пулом можуть бути релевантні документи. Повнота на глибині 100 визначається відносно відомих релевантних документів пулу; зарахування неоцінених документів до нерелевантних може сприяти системам, які формували пул, і не вимірює повноту за всім корпусом. Відсутність статистично значущої різниці не доводить еквівалентності.

\subsection{Польове дослідження одного автора}
Польове дослідження (розділ~\ref{sec:field}) використовує архів однієї людини та набори запитів, які автор писав під час розроблення системи, зокрема для її перевірки. Це не випадкова вибірка реальних інформаційних потреб, а систему налаштовували на тому самому архіві. Результати характеризують поведінку в цьому середовищі, а не узагальнювану ефективність.

\subsection{Вимірювання}
Усі вимірювання витрат стосуються одного ноутбука у звичайному сеансі з іншими програмами; навантаження перед кожним кроком записується, але не контролюється. Таймер запиту охоплює кодування й пошук фрагментів, але не згортання до документів, визначення шляхів і відображення інтерфейсу. Лічильники ресурсів процесу охоплюють увесь виміряний прохід, включно зі згортанням, тому їхня область вимірювання відрізняється від таймера запиту. Індексування на CPU виміряно на п'ятій частині вибірки GPU. Експеримент масштабування використовує синтетичні вектори, утворені зі справжніх; повнота HNSW на них може відрізнятися від повноти на реальній колекції такого розміру.

\subsection{Валідність вимірюваних понять}
Судді бачать назву й перші 1\,200 символів документа, а не весь документ. Тому релевантність, що проявляється лише далі, може бути недооцінена; однаковий фрагмент показується для всіх систем, проте це сприяє документам із релевантністю на початку. Оцінки 0/1/2 стискають судження, які люди можуть робити інакше: модельне «частково релевантний» не дорівнює людському. Основна метрика nDCG@10 за мітками пулу винагороджує високе ранжування оцінених релевантних документів; вона не вимірює, чи людина швидше знаходить те, що пам'ятає.

\subsection{Завдання зіставлення джерельних пар}
Завдання із зображеннями й мовленням визначають пов'язаний джерельний об'єкт або речення. Ці пари не дають вичерпних оцінок кожного семантично спорідненого кандидата. Підписи й паралельні транскрипти також відрізняються від незалежно перефразованих запитів до архіву. Для поданих точкових оцінок довірчі інтервали не обчислювалися.
